\documentclass[10pt]{article}

\usepackage[utf8]{inputenc}

\usepackage[T1]{fontenc}

\usepackage{amsmath}

\usepackage{amsfonts}

\usepackage{amssymb}

\usepackage[version=4]{mhchem}

\usepackage{stmaryrd}

\usepackage{bbold}

\usepackage{graphicx}

\usepackage[export]{adjustbox}

\graphicspath{ {./images/} }



\begin{document}

где $H: M \times \Delta \rightarrow \mathbb{R}$ - функция Гамильтона, а $\Delta \subset \mathbb{R}$ - интервал времени. Экстремалями этого функционала являются траектории гамильтоновой системы с гамильтонианом $H$.



См. также Симплектическая геометрия. М. Э. Казарян. ДЕЙСТВИЕ по Гам и льтону - см. Вариационные принципы классической механики, Гамильтона - Остроградского принцип.



ДЕЙСТВИЕ по Лагранжу - см. Вариационные принципы классической механики, Гамильтона - Якоби уравнение, Лагранжа принцип.



ДЕЙСТВИЕ по Якоби-см. Вариационные принципы классической механики, Якоби принцип.



ДЕЙСТВИЕ-УГОЛ ПЕРЕМЕННЫЕ - см. Гамильтонова система.



ДЕЙСТВИЯ ФУНКЦИОНАЛ - см. Случайное возмущение. ДЕКАРТОВО ПРОИЗВЕДЕНИЕ - см. Прямое произведение.



ДЕЛОНЕ МЕТОД - см. Усреднения метод.



ДЕЛЬТА-ПОСЛЕДОВАТЕЛЬНОСТЬ - см. РегУляризация обобщенной функции.



ДЕЛЬТА-ФУнКЦИЯ, $\delta(x)$,- то же, что Дирака дельта-функция. См. также Обобщенная функция.



ДЕСИТТЕРА ГPУПกА - см. Дe Cuтmepa груnпа.



ДЕТАЛЬНОГО РАВНОВЕСИЯ ПРИНЦИП - см. Основное кинетическое уравнение.



ДЕТЕРМИНИРОВАННОГО ХАОСА ТЕОРИЯ - см. Эргодическая теорема.



ДЕТОНАЦИИ ВОЛНЫ - см. УдарнЫе волнЫ.



ДЕФЕКТ МАСС - см. Macca.



ДЕФОРМАЦИЙ ТЕНЗОР - симметрический тензор, задающий геометрию деформированного твердого тела. Д.т. выражается формулой:



\[

\varepsilon_{i k}=\frac{1}{2}\left(\frac{\partial u_{i}}{\partial x_{k}}+\frac{\partial u_{k}}{\partial x_{i}}+\frac{\partial u_{i}}{\partial x_{i}} \frac{\partial u_{k}}{\partial x_{k}}\right), \quad i, k=1,2,3 .

\]



Здесь $u_{i}=x_{i}^{\prime}-x_{i}-$ компоненты вектора смещений, где $x_{i}-$ координата точки тела до смешения, $x_{i}^{\prime}$ - после смещения. Изменение элемента длины $\boldsymbol{d l}^{\prime}$ деформированного тела связано с Д. т. соотношением



\[

d l^{2}=\left(\delta_{i k}+2 \varepsilon_{i k}\right) d x_{i} d x_{k},

\]



где $\delta_{i k} d x_{i} d x_{k}$ - квадрат элемента длины недеформированного тела, $\delta_{i k}$ - символ Кронекера. ДЕФОРМАЦИЯ - см. Бифуркация, Версальная деформация. ДЖАФФЕ КЛАСС - класс пространств типа $S^{\prime}$, основные пространства которых содержат разложение единицы. Д. к. введен в середине 60-х гг. 20 в. А. Джаффе [2] в связи с проблемой локальности взаимодействия в аксиоматич. квантовой теории поля.



Лum:: [1] Боголюбов Н.Н., Логунов А.А., Тодоров И. Т., Основы аксиоматического подхода в квантовой теории поля, M., 1969, c. 293; [2] J affe A., «Phys. Rev.», 1967, v. 158, p. 1454-61.



ДХET - см. Расслоение струй.



\includegraphics[max width=\textwidth, center]{2023_07_08_3c43317f01c89fc83b46g-1}

см. Вейеритрасса элиитииеские функции.



ДИАГОНАЛЬНАЯ АППРОКСИМАЦИЯ Паде - см. Паде аппроксимация.



ДИАГРАММ МЕТОД в те ор и и турбуле нтности метод, основанный на графическом представлении в виде диаграмм Фейнмана отдельных членов разложения решений уравнений гидродинамики в формальный ряд теории возмущений по степеням константы взаимодействия - числу Рейнольдса. Свойства диаграмм позволяют производить частичное суммирование ряда с целью получения замкнутых интегральных уравнений, содержащих в качестве неизвестных искомые величины турбулентного потока.



Лит.: [1] М он и н А. С., Я гл о м А. М., Статистическая гидромеханика, ч. 2, М., 1967; [2] Гледзе р Е. Б., Мон и н А. С., «Успехи $\begin{array}{lll}\text { математич. наук॰, 1974, т. 29, в. 3, с. 111-59. } & \text { Е. Б. Гледзер. }\end{array}$ ДИАНАЛИТИЧЕСКАЯ КАРТА - см. Клейнова поверхHость.



ДИАНАЛИТИЧЕСКИЙ АТЛАС - см. Клейнова повер $x-$ ность.



ДИВЕРГЕНЦИЯ ве ктор н о го п о ля $\boldsymbol{a}$ (в точке $x, y, z)$ скалярная величина



\[

\operatorname{div} \boldsymbol{a}=\frac{\partial a_{x}}{\partial x}+\frac{\partial a_{y}}{\partial y}+\frac{\partial a_{z}}{\partial z}

\]



где $a_{x}, a_{y}, a_{z}-$ координаты вектора $\boldsymbol{a}$. Если вектор $\boldsymbol{a}-$ скорость частицы в установившемся течении несжимаемой жидкости, то div $\boldsymbol{a}$ в точке с координатами $x, y, z$ обозначает интенсивность источника (div $\boldsymbol{a}>0$ ) или стока $(\operatorname{div} a<0)$, находяшегося в этой точке, или отсутствие источника или стока $(\operatorname{div} a=0)$. Д. есть предел отношения потока векторного поля через замкнутую поверхность, окружающую данную точку, к объему, ограничиваемому ею, когда эта поверхность стягивается к точке.



Свойства Д.:



\[

\begin{gathered}

\operatorname{div}(\boldsymbol{a}+\boldsymbol{b})=\operatorname{div} \boldsymbol{a}+\operatorname{div} \boldsymbol{b}, \\

\operatorname{div}(\varphi \boldsymbol{a})=\varphi \operatorname{div} \boldsymbol{a}+\boldsymbol{a} \operatorname{grad} \varphi, \\

\operatorname{div} \operatorname{rot} \boldsymbol{a}=0, \\

\operatorname{div} \operatorname{grad} \boldsymbol{a}=\Delta \varphi, \Delta-\text { оператор Лапласа, } \\

\operatorname{div}[\boldsymbol{a}, \boldsymbol{b}]=(\boldsymbol{b}, \operatorname{rot} \boldsymbol{a})-(\boldsymbol{a}, \operatorname{rot} \boldsymbol{b}) .

\end{gathered}

\]



См. также Векторный анализ. ДиЗъЮНКТНОСТЬ - отношение в классе унитарных представлений. А именно, унитарные представления $T_{1}$ и $T_{2}$ называются дизъюн ктным и, если не существует ненулевого сплетающего оператора из $T_{1}$ в $T_{2}$. Представления $T_{1}$ и $T_{2}$ дизъюнктны тогда и только тогда, когда никакое подпредставление представления $T_{1}$ не эквивалентно никакому подпредставлению представления $T_{2}$. Ю. А. Неретин. ДИКВАРК - см. Лептокварк.



ДИККЕ МОДЕЛЬ - модель, описывающая взаимодействие системы $N$ двухуровневых излучателей с квантованным электромагнитным полем. В общем случае Д.м. задается гамильтонианом



\[

\begin{gathered}

\boldsymbol{H}=\frac{1}{2} \sum_{i=1}^{N} \hbar \varepsilon_{i} \sigma_{i}^{z}+\sum_{k} \hbar \omega_{k} a_{k}^{+} a_{k}- \\

-\sum_{i=1}^{N} \sum_{k}\left(g_{i k} \sigma_{i}^{+}+g_{i k}^{*} \sigma_{i}^{-}\right)\left(a_{k} e^{i \overrightarrow{\boldsymbol{k}} \vec{r}_{i}}+a_{k}^{+} e^{-i \overrightarrow{\boldsymbol{k}} \vec{r}_{i}}\right)

\end{gathered}

\]



в к-ром $\hbar \varepsilon_{i}-$ энергия межуровневого перехода $i$-го излучателя, $\hbar \omega_{k}-$ энергия кванта излучения с волновым вектором $\boldsymbol{k}$ и поляризацией $v$, суммирование по $k \equiv\{\boldsymbol{k}, \boldsymbol{v}\}$ обозначает суммирование и по $k$, и по $v, a_{k}^{+}$и $a_{k}$-операторы рождения и гибели фотона, $\sigma_{i}^{ \pm} \equiv \frac{1}{2}\left(\sigma_{i}^{x} i \sigma_{i}^{y}\right)-$ лестничные операторы, $\sigma_{i}^{\alpha}, \alpha=x, y, z$,- матрицы Паули, ассоциированные с $i$-м излучателем, находящимся в точке $\boldsymbol{r}_{i}$,



\[

g_{i k}=\left(2 \pi \hbar / V \omega_{k}\right)^{1 / 2}\left(\boldsymbol{d}_{i}, \boldsymbol{e}_{\mathrm{v}}\right),

\]



$V$ - объем квантования, $d_{i}-$ дипольный момент перехода $i$-го излучателя, $\boldsymbol{e}_{v}-$ единичный вектор поляризации с номером $v$.



Обычно для вычислений используется Д. м. в приближении вращающейся волны, когда в гамильтониане взаимодействия оставляют лишь резонансные слагаемые типа $\sigma_{i}^{+} a_{k}$ и $\sigma_{i}^{-} a_{k}^{+}$, а нерезонансные вида $\sigma_{i}^{+} a_{k}^{+}$и $\sigma_{i}^{-} a_{k}$ отбрасывают. Если длина волны излучения гораздо больше характерных размеров системы, то переходят к сосредоточенной Д. м., гамильтониан к-рой не содержит суммирований по $k$. Это эквивалентно переходу к одномодовой Д. м., в приближении вращающейся волны, задающейся гамильтонианом



\[

H=\frac{1}{2} \sum_{i=1}^{N} \hbar \varepsilon_{i} \sigma_{i}^{z}+\hbar \omega a^{+} a-\sum_{i=1}^{N}\left(g_{i} \sigma_{i}^{+} a+g_{i}^{*} \sigma_{i}^{-} a^{+}\right) .

\]





\end{document}